\documentclass[10pt,a4paper]{amsart}
\usepackage{amsmath,amssymb,amsthm,graphicx}
% For commutative diagrams we use Xy-pic
\usepackage[all]{xy}

% We use 2cell for 2-commutative diagrams.
\xyoption{2cell}
% Theorem environments.
%
\theoremstyle{plain}
\newtheorem{theorem}[subsection]{Theorem}
\newtheorem{proposition}[subsection]{Proposition}
\newtheorem{lemma}[subsection]{Lemma}

\theoremstyle{definition}
\newtheorem{definition}[subsection]{Definition}
\newtheorem{example}[subsection]{Example}
\newtheorem{exercise}[subsection]{Exercise}
\newtheorem{situation}[subsection]{Situation}

\theoremstyle{remark}
\newtheorem{remark}[subsection]{Remark}
\newtheorem{remarks}[subsection]{Remarks}

\numberwithin{equation}{subsection}

% Macros
%
\def\lim{\mathop{\mathrm{lim}}\nolimits}
\def\colim{\mathop{\mathrm{colim}}\nolimits}
\def\Spec{\mathop{\mathrm{Spec}}}
\def\Hom{\mathop{\mathrm{Hom}}\nolimits}
\def\Ext{\mathop{\mathrm{Ext}}\nolimits}
\def\SheafHom{\mathop{\mathcal{H}\!\mathit{om}}\nolimits}
\def\SheafExt{\mathop{\mathcal{E}\!\mathit{xt}}\nolimits}
\def\Sch{\mathit{Sch}}
\def\Mor{\mathop{\mathrm{Mor}}\nolimits}
\def\Ob{\mathop{\mathrm{Ob}}\nolimits}
\def\Sh{\mathop{\mathit{Sh}}\nolimits}
\def\NL{\mathop{N\!L}\nolimits}
\def\CH{\mathop{\mathrm{CH}}\nolimits}
\def\proetale{{pro\text{-}\acute{e}tale}}
\def\etale{{\acute{e}tale}}
\def\QCoh{\mathit{QCoh}}
\def\Ker{\mathop{\mathrm{Ker}}}
\def\Im{\mathop{\mathrm{Im}}}
\def\Coker{\mathop{\mathrm{Coker}}}
\def\Coim{\mathop{\mathrm{Coim}}}

% Boxtimes
%
\DeclareMathSymbol{\boxtimes}{\mathbin}{AMSa}{"02}

%
% Macros for moduli stacks/spaces
%
\def\QCohstack{\mathcal{QC}\!\mathit{oh}}
\def\Cohstack{\mathcal{C}\!\mathit{oh}}
\def\Spacesstack{\mathcal{S}\!\mathit{paces}}
\def\Quotfunctor{\mathrm{Quot}}
\def\Hilbfunctor{\mathrm{Hilb}}
\def\Curvesstack{\mathcal{C}\!\mathit{urves}}
\def\Polarizedstack{\mathcal{P}\!\mathit{olarized}}
\def\Complexesstack{\mathcal{C}\!\mathit{omplexes}}
% \Pic is the operator that assigns to X its picard group, usage \Pic(X)
% \Picardstack_{X/B} denotes the Picard stack of X over B
% \Picardfunctor_{X/B} denotes the Picard functor of X over B
\def\Pic{\mathop{\mathrm{Pic}}\nolimits}
\def\Picardstack{\mathcal{P}\!\mathit{ic}}
\def\Picardfunctor{\mathrm{Pic}}
\def\Deformationcategory{\mathcal{D}\!\mathit{ef}}
\usepackage{cleveref}
\crefname{originallemma}{original lemma}{original lemmas}
\Crefname{originallemma}{Original lemma}{Original lemmas}
\makeatletter
\newlabel{categories-definition-exact}{{0034}{source}{}{sourcepointer-1}{}}
\newlabel{categories-definition-exact@cref}{{[originallemma][0][]0034}{[1][0][]source}{}{}{}}
\makeatother
\crefname{originallemma}{original lemma}{original lemmas}
\Crefname{originallemma}{Original lemma}{Original lemmas}
\makeatletter
\newlabel{sites-proposition-get-morphism}{{00X6}{source}{}{sourcepointer-2}{}}
\newlabel{sites-proposition-get-morphism@cref}{{[originallemma][0][]00X6}{[1][0][]source}{}{}{}}
\makeatother
\crefname{originalequation}{original equation}{original equations}
\Crefname{originalequation}{Original equation}{Original equations}
\makeatletter
\newlabel{equation-functorial}{{04I1}{source}{}{sourcepointer-3}{}}
\newlabel{equation-functorial@cref}{{[originalequation][0][]04I1}{[1][0][]source}{}{}{}}
\makeatother
\crefname{originallemma}{original lemma}{original lemmas}
\Crefname{originallemma}{Original lemma}{Original lemmas}
\makeatletter
\newlabel{sites-section-sheaves-algebraic-structures}{{00YR}{source}{}{sourcepointer-4}{}}
\newlabel{sites-section-sheaves-algebraic-structures@cref}{{[originallemma][0][]00YR}{[1][0][]source}{}{}{}}
\makeatother
\crefname{originallemma}{original lemma}{original lemmas}
\Crefname{originallemma}{Original lemma}{Original lemmas}
\makeatletter
\newlabel{sites-modules-lemma-flat-pullback-exact}{{04JC}{source}{}{sourcepointer-5}{}}
\newlabel{sites-modules-lemma-flat-pullback-exact@cref}{{[originallemma][0][]04JC}{[1][0][]source}{}{}{}}
\makeatother
\crefname{originaltheorem}{original theorem}{original theorems}
\Crefname{originaltheorem}{Original theorem}{Original theorems}
\makeatletter
\newlabel{theorem-exactness-stalks}{{03PU}{source}{}{sourcepointer-6}{}}
\newlabel{theorem-exactness-stalks@cref}{{[originaltheorem][0][]03PU}{[1][0][]source}{}{}{}}
\makeatother
\crefname{originallemma}{original lemma}{original lemmas}
\Crefname{originallemma}{Original lemma}{Original lemmas}
\makeatletter
\newlabel{sites-lemma-point-morphism-sites}{{05V1}{source}{}{sourcepointer-7}{}}
\newlabel{sites-lemma-point-morphism-sites@cref}{{[originallemma][0][]05V1}{[1][0][]source}{}{}{}}
\makeatother
\crefname{originallemma}{original lemma}{original lemmas}
\Crefname{originallemma}{Original lemma}{Original lemmas}
\makeatletter
\newlabel{lemma-stalk-exact}{{03PT}{source}{}{sourcepointer-8}{}}
\newlabel{lemma-stalk-exact@cref}{{[originallemma][0][]03PT}{[1][0][]source}{}{}{}}
\makeatother
\crefname{originalequation}{original equation}{original equations}
\Crefname{originalequation}{Original equation}{Original equations}
\makeatletter
\newlabel{equation-pullback}{{04I3}{source}{}{sourcepointer-9}{}}
\newlabel{equation-pullback@cref}{{[originalequation][0][]04I3}{[1][0][]source}{}{}{}}
\makeatother
\crefname{originallemma}{original lemma}{original lemmas}
\Crefname{originallemma}{Original lemma}{Original lemmas}
\makeatletter
\newlabel{categories-lemma-yoneda}{{001P}{source}{}{sourcepointer-10}{}}
\newlabel{categories-lemma-yoneda@cref}{{[originallemma][0][]001P}{[1][0][]source}{}{}{}}
\makeatother
\crefname{originallemma}{original lemma}{original lemmas}
\Crefname{originallemma}{Original lemma}{Original lemmas}
\makeatletter
\newlabel{sites-lemma-pullback-representable-sheaf}{{04D3}{source}{}{sourcepointer-11}{}}
\newlabel{sites-lemma-pullback-representable-sheaf@cref}{{[originallemma][0][]04D3}{[1][0][]source}{}{}{}}
\makeatother
\title{Exactness of the small inverse image and its effect on stalks and representable sheaves}
\author{The Stacks Project Authors}
\date{stacks-project:etale-cohomology.tex:local-snapshot-20261011; GNU FDL 1.2 or later}
\begin{document}
\maketitle
\noindent\textit{Editorial note.} All delivered mathematical text is the author's own, copied from the pinned Stacks Project chapter file etale-cohomology.tex; no statement or proof step was rewritten, reordered or removed.\par\medskip
\noindent\textit{Editorial note.} The proof is the author's complete argument for all four parts; in part (4) the author gives both the direct computation and the alternative route through the previous parts.\par\medskip
\noindent\textit{Editorial note.} Every labelled reference occurring inside the counted spans is carried in assembly.external\_references as a non-load-bearing navigation pointer (same-chapter ones at their definition site, cross-chapter ones at the invocation inside the counted span), so no load-bearing mathematics is replaced by a stub. No bibliography citation occurs inside the counted spans.\par\medskip
\noindent\textit{Editorial note.} The checked chapter list for same-name collisions is the set of stacks chapters staged in this repository: chow, cohomology, divisors, etale-cohomology, flat, more-algebra, more-morphisms, morphisms, spaces-more-morphisms, varieties together with this card's chapter etale-cohomology.\par\medskip
\noindent\textit{Editorial note.} The preamble is the chapter's own native preamble excerpt already used for this source: the Xy-pic setup (native preamble lines 14-18) and the native theorem declarations with the mathematical macros (native preamble lines 154-217).\par\medskip
\noindent\textit{Allowed standard prerequisites (editorial).}\begin{itemize}
\item Small inverse image f\textasciicircum{}\{-1\} on sheaves on the small etale site. Same chapter, the functor f\textasciicircum{}\{-1\} attached to a morphism of schemes and its description by colimits over etale neighbourhoods. Supplies the functors and identifications appearing in all four parts of the statement.
\item Stalks of a sheaf at a geometric point. Same chapter, the stalk functor at a geometric point of the small etale site, and its behaviour under pullback. Supplies parts (2) and (3), including the identification of the stalk of the pullback with the stalk at the composite geometric point.
\item Sheaf represented by an etale scheme and its sections. Same chapter, the representable sheaf h\_V on Y\_etale for V -> Y etale and the computation of f\textasciicircum{}\{-1\}h\_V. Supplies part (4), namely f\textasciicircum{}\{-1\}h\_V = h\_\{X x\_Y V\}.
\end{itemize}\medskip
\ifdefined\hypertarget\hypertarget{sourcepointer-1}{}\fi
\noindent\textit{Source reference.} Original lemma 0034: sources/stacks-etale-cohomology/pinned-original-source.tex, source lines 4658--4658. Navigation pointer invoked at proof line 4658; the counted span contains the pointer only and the cited result is not reproduced inside it, so the pointer is registered as a non-load-bearing navigation stub in the established form for Stacks material. No bibliography citation is involved.\par\medskip
\ifdefined\hypertarget\hypertarget{sourcepointer-2}{}\fi
\noindent\textit{Source reference.} Original lemma 00X6: sources/stacks-etale-cohomology/pinned-original-source.tex, source lines 4640--4640. Navigation pointer invoked at proof line 4640; the counted span contains the pointer only and the cited result is not reproduced inside it, so the pointer is registered as a non-load-bearing navigation stub in the established form for Stacks material. No bibliography citation is involved.\par\medskip
\ifdefined\hypertarget\hypertarget{sourcepointer-3}{}\fi
\noindent\textit{Source reference.} Original equation 04I1: sources/stacks-etale-cohomology/pinned-original-source.tex, source lines 4432--4432. Navigation pointer invoked at proof line 4641; the counted span contains the pointer only and the cited result is not reproduced inside it, so the pointer is registered as a non-load-bearing navigation stub in the established form for Stacks material. No bibliography citation is involved.\par\medskip
\ifdefined\hypertarget\hypertarget{sourcepointer-4}{}\fi
\noindent\textit{Source reference.} Original lemma 00YR: sources/stacks-etale-cohomology/pinned-original-source.tex, source lines 4680--4680. Navigation pointer invoked at proof line 4680; the counted span contains the pointer only and the cited result is not reproduced inside it, so the pointer is registered as a non-load-bearing navigation stub in the established form for Stacks material. No bibliography citation is involved.\par\medskip
\ifdefined\hypertarget\hypertarget{sourcepointer-5}{}\fi
\noindent\textit{Source reference.} Original lemma 04JC: sources/stacks-etale-cohomology/pinned-original-source.tex, source lines 4682--4682. Navigation pointer invoked at proof line 4682; the counted span contains the pointer only and the cited result is not reproduced inside it, so the pointer is registered as a non-load-bearing navigation stub in the established form for Stacks material. No bibliography citation is involved.\par\medskip
\ifdefined\hypertarget\hypertarget{sourcepointer-6}{}\fi
\noindent\textit{Source reference.} Original theorem 03PU: sources/stacks-etale-cohomology/pinned-original-source.tex, source lines 3546--3546. Navigation pointer invoked at proof line 4684; the counted span contains the pointer only and the cited result is not reproduced inside it, so the pointer is registered as a non-load-bearing navigation stub in the established form for Stacks material. No bibliography citation is involved.\par\medskip
\ifdefined\hypertarget\hypertarget{sourcepointer-7}{}\fi
\noindent\textit{Source reference.} Original lemma 05V1: sources/stacks-etale-cohomology/pinned-original-source.tex, source lines 4688--4688. Navigation pointer invoked at proof line 4688; the counted span contains the pointer only and the cited result is not reproduced inside it, so the pointer is registered as a non-load-bearing navigation stub in the established form for Stacks material. No bibliography citation is involved.\par\medskip
\ifdefined\hypertarget\hypertarget{sourcepointer-8}{}\fi
\noindent\textit{Source reference.} Original lemma 03PT: sources/stacks-etale-cohomology/pinned-original-source.tex, source lines 3456--3456. Navigation pointer invoked at proof line 4690; the counted span contains the pointer only and the cited result is not reproduced inside it, so the pointer is registered as a non-load-bearing navigation stub in the established form for Stacks material. No bibliography citation is involved.\par\medskip
\ifdefined\hypertarget\hypertarget{sourcepointer-9}{}\fi
\noindent\textit{Source reference.} Original equation 04I3: sources/stacks-etale-cohomology/pinned-original-source.tex, source lines 4632--4632. Navigation pointer invoked at proof line 4697; the counted span contains the pointer only and the cited result is not reproduced inside it, so the pointer is registered as a non-load-bearing navigation stub in the established form for Stacks material. No bibliography citation is involved.\par\medskip
\ifdefined\hypertarget\hypertarget{sourcepointer-10}{}\fi
\noindent\textit{Source reference.} Original lemma 001P: sources/stacks-etale-cohomology/pinned-original-source.tex, source lines 4713--4713. Navigation pointer invoked at proof line 4713; the counted span contains the pointer only and the cited result is not reproduced inside it, so the pointer is registered as a non-load-bearing navigation stub in the established form for Stacks material. No bibliography citation is involved.\par\medskip
\ifdefined\hypertarget\hypertarget{sourcepointer-11}{}\fi
\noindent\textit{Source reference.} Original lemma 04D3: sources/stacks-etale-cohomology/pinned-original-source.tex, source lines 4721--4721. Navigation pointer invoked at proof line 4721; the counted span contains the pointer only and the cited result is not reproduced inside it, so the pointer is registered as a non-load-bearing navigation stub in the established form for Stacks material. No bibliography citation is involved.\par\medskip
\begin{lemma}
\label{lemma-stalk-pullback}
Let $f : X \to Y$ be a morphism of schemes.
\begin{enumerate}
\item The functor
$f^{-1} : \textit{Ab}(Y_\etale) \to \textit{Ab}(X_\etale)$
is exact.
\item The functor
$f^{-1} : \Sh(Y_\etale) \to \Sh(X_\etale)$
is exact, i.e., it commutes with finite limits and colimits, see
Categories, Definition \ref{categories-definition-exact}.
\item Let $\overline{x} \to X$ be a geometric point.
Let $\mathcal{G}$ be a sheaf on $Y_\etale$.
Then there is a canonical identification
$$
(f^{-1}\mathcal{G})_{\overline{x}} = \mathcal{G}_{\overline{y}}.
$$
where $\overline{y} = f \circ \overline{x}$.
\item For any $V \to Y$ \'etale we have $f^{-1}h_V = h_{X \times_Y V}$.
\end{enumerate}
\end{lemma}

\begin{proof}
The exactness of $f^{-1}$ on sheaves of sets is a consequence of
Sites, Proposition \ref{sites-proposition-get-morphism}
applied to our functor $u$ of Equation (\ref{equation-functorial}).
In fact the exactness of pullback is part of the definition of
a morphism of topoi (or sites if you like). Thus we see (2) holds.
It implies part (1) since given an abelian sheaf $\mathcal{G}$ on
$Y_\etale$
the underlying sheaf of sets of $f^{-1}\mathcal{F}$ is the same
as $f^{-1}$ of the underlying sheaf of sets of $\mathcal{F}$, see
Sites, Section \ref{sites-section-sheaves-algebraic-structures}.
See also
Modules on Sites, Lemma \ref{sites-modules-lemma-flat-pullback-exact}.
In the literature (1) and (2) are sometimes deduced from (3) via
Theorem \ref{theorem-exactness-stalks}.

\medskip\noindent
Part (3) is a general fact about stalks of pullbacks, see
Sites, Lemma \ref{sites-lemma-point-morphism-sites}.
We will also prove (3) directly as follows. Note that by
Lemma \ref{lemma-stalk-exact}
taking stalks commutes with sheafification.
Now recall that $f^{-1}\mathcal{G}$ is the sheaf
associated to the presheaf
$$
U \longrightarrow \colim_{U \to X \times_Y V} \mathcal{G}(V),
$$
see Equation (\ref{equation-pullback}).
Thus we have
\begin{align*}
(f^{-1}\mathcal{G})_{\overline{x}}
& = \colim_{(U, \overline{u})} f^{-1}\mathcal{G}(U) \\
& = \colim_{(U, \overline{u})}
\colim_{a : U \to X \times_Y V} \mathcal{G}(V) \\
& = \colim_{(V, \overline{v})} \mathcal{G}(V) \\
& = \mathcal{G}_{\overline{y}}
\end{align*}
in the third equality the pair $(U, \overline{u})$ and the map
$a : U \to X \times_Y V$ corresponds to the pair $(V, a \circ \overline{u})$.

\medskip\noindent
Part (4) can be proved in a similar manner by identifying the colimits
which define $f^{-1}h_V$. Or you can use
Yoneda's lemma (Categories, Lemma \ref{categories-lemma-yoneda})
and the functorial equalities
$$
\Mor_{\Sh(X_\etale)}(f^{-1}h_V, \mathcal{F}) =
\Mor_{\Sh(Y_\etale)}(h_V, f_*\mathcal{F}) =
f_*\mathcal{F}(V) = \mathcal{F}(X \times_Y V)
$$
combined with the fact that representable presheaves are sheaves. See also
Sites, Lemma \ref{sites-lemma-pullback-representable-sheaf}
for a completely general result.
\end{proof}

\end{document}
